\labsection{6}{Threads, Banker's algorithm, and deadlock reasoning}{sec:lab06}
\begin{labproject}{Combine thread lifecycle with resource-allocation safety analysis}
\textbf{Coverage.} Chapter 6.\\
\textbf{Source files.} Three C programs in \texttt{Lab06}.\\
\textbf{Goal.} Practice POSIX thread creation/joining and matrix-based resource reasoning.

\medskip\noindent\textbf{Task A --- Banker's safety algorithm.}\par
After input, calculate every Need entry with \texttt{Need = Max - Allocation}. Maintain \texttt{Work} and \texttt{Finish}. Record the process selected on each pass and the Work vector after its resources are released.

\medskip\noindent\textbf{Task B --- Detection terminology.}\par
Compare \texttt{Task01.c} and \texttt{Task02.c}. Both derive Need from Max. Explain why a strict deadlock-detection algorithm is conceptually different from a Banker's safe-state test.

\medskip\noindent\textbf{Task C --- POSIX threads.}\par
Create two threads, identify them with \texttt{pthread\_self()}, compare identifiers with \texttt{pthread\_equal()}, and join both before process exit.

\begin{tryit}
Remove the busy delay and run the thread program several times. Then add a shared counter without a mutex, followed by a mutex-protected version. Compare results.
\end{tryit}
\end{labproject}

\begin{labdeliverable}
Submit one safe and one unsafe Banker's input with full Work/Finish traces, a paragraph distinguishing unsafe from deadlocked, and thread output showing successful creation and joining.
\end{labdeliverable}
